\documentclass[12pt]{article}

\usepackage{amsmath}
\usepackage{amssymb}
\usepackage{amsfonts}
\usepackage{geometry}
\usepackage{setspace}
\usepackage{booktabs}
\usepackage[T1]{fontenc}
\usepackage[utf8]{inputenc}

\geometry{margin=1in}
\setstretch{1.2}

\begin{document}

\subsection*{Carr and Wu (2003): What Type of Process Underlies Options? A Simple Robust Test}

\subsubsection*{Introduction}

A fundamental question in asset pricing is whether stock prices are best described
by a purely continuous process (PC), a pure jump process (PJ), or a combination
of both (CJ). The direct approach --- examining time-series data --- is unreliable
unless the sampling frequency is extremely high, and high-frequency data
introduce microstructure noise. Carr and Wu (2003) sidestep this problem
entirely by looking at option prices instead. Their key insight is that, although
discretely sampled paths from the three process types look nearly identical, they
have dramatically different implications for short-maturity option prices.

\subsubsection*{The Core Method: Term Decay Plots}

Under a risk-neutral measure, the asset price dynamics are modelled as:

\begin{equation}
    \frac{dS_t}{S_{t-}} = (r-q)\,dt + \sigma_t\,dW_t
    + \int_{\mathbb{R}^0}(e^x - 1)\bigl[\mu(dx,dt) - \nu_t(x)\,dx\,dt\bigr]
\end{equation}

where $W_t$ is a standard Brownian motion, $\mu(dx,dt)$ counts jumps of log-size
$x$, and $\nu_t(x)$ is the compensator (local jump density). The first stochastic
term is the purely continuous martingale component; the integral is the purely
discontinuous (jump) martingale component.

The paper's main theoretical result decomposes the time value of a European
option into two parts --- one driven by the continuous component and one by the
jump component:

\begin{equation}
    TV_0(K,T) = e^{-rT}\int_0^T
    \Bigl[\tfrac{1}{2}q(K,t)K^2\,\mathbb{E}_0[\sigma_t^2\,|\,F_{t-}=K]
    + \mathbb{E}_0[F_{t-}\,\nu_t^0(k)]\Bigr]\,dt
\end{equation}

where $q(K,t)$ is the risk-neutral density of the forward price, $k=\ln(K/F_{t-})$
is log-moneyness, and $\nu_t^0(k)$ is the double tail of the jump compensator.
As $T\downarrow 0$, this decomposition reveals that the speed at which option
prices converge to zero depends entirely on which components are present, as
summarised in Table~\ref{tab:decay}.

\begin{table}[h]
\centering
\caption{Asymptotic Decay Rates of Short-Maturity Option Prices}
\label{tab:decay}
\begin{tabular}{lll}
\toprule
Process Type & OTM Options & ATM Options \\
\midrule
Purely Continuous (PC)  & $O(e^{-c/T}),\;c>0$ & $O(\sqrt{T})$ \\
Pure Jump (PJ)          & $O(T)$               & $O(T^p),\;p\in(0,1]$ \\
Combined (CJ)           & $O(T)$               & $O(T^p),\;p\in(0,\tfrac{1}{2}]$ \\
\bottomrule
\end{tabular}
\end{table}

The intuition is straightforward. Under a purely continuous process, an
out-of-the-money option can only move in the money through gradual
diffusion, which is exponentially unlikely over a very short horizon. When
jumps are present, however, the stock can leap across the strike in a single
instant, so OTM option prices decay only linearly in $T$. For ATM options,
the $O(\sqrt{T})$ rate signals an infinite-variation component (either diffusion
or an infinite-variation jump process), while $O(T)$ signals a pure finite-variation
jump process.

To apply the test empirically, the authors plot $\ln(P/T)$ against $\ln(T)$,
calling this a \textit{term decay plot}. A flat slope near zero for OTM options
signals a jump component; a strongly positive, concave curve signals a
purely continuous process. For ATM options, a straight line with slope
near $-0.5$ signals an infinite-variation component; a flat line signals
finite-variation jumps only.

\subsubsection*{Empirical Findings: S\&P 500 Index Options}

Carr and Wu apply the method to daily S\&P 500 index option data from
April 1999 to May 2000 (273 business days). They fit a second-order polynomial,

\begin{equation}
    \ln(P/T) = a(\ln T)^2 + b(\ln T) + c,
\end{equation}

to the term decay plot at each date and moneyness level, extracting slope
and curvature estimates. The results clearly support a combined
continuous-jump (CJ) process. OTM term decay plots are mostly flat at
short maturities (slopes near zero), confirming the presence of a jump
component. ATM term decay plots are consistently negatively sloped and
approximately linear with slopes averaging $-0.413$, close to the
$-0.5$ theoretical value for a continuous component, confirming that
an infinite-variation component is always present. Importantly, the jump
component varies substantially over time --- on some days it is almost
undetectable --- while the continuous component is constantly and
stably felt throughout the sample.

\subsubsection*{Strengths, Limitations, and Relevance to QQQ Options}

The main strength of the method is its robustness: it does not require a
parametric model and is not confined to a Markovian or single-factor
framework. It can detect the type of process directly from observable option
prices, without needing high-frequency return data. The simulation evidence
confirms that the asymptotic behavior is visible for options maturing within
20 days, which are readily available in liquid markets.

The main limitation is that the test cannot always distinguish between a
continuous martingale and an infinite-variation pure jump process, since
both generate a similar $O(\sqrt{T})$ ATM decay rate. The method also
operates under the assumption that the compensator $\sigma_t$ and $\nu_t$
are bounded near $t=0$; deterministically explosive volatility could
mimic jump-like decay even in a purely continuous model.

For our project on QQQ options, the Carr and Wu framework provides
direct empirical motivation for including a jump component in the pricing
model. The NASDAQ-100's technology-heavy composition makes it prone
to sudden, large price moves, and if QQQ option prices exhibit the flat
short-maturity OTM decay pattern documented here, this constitutes
market-based evidence that jumps are priced into the index. The finding
that the continuous component is always present further supports the use
of combined jump-diffusion models --- such as Merton (1976) or Kou (2002)
--- rather than pure jump specifications. The time-variation in the jump
component also suggests that models with stochastic jump intensity may
be necessary for capturing the full dynamics of the QQQ implied volatility
surface across different market regimes.

\end{document}
